\documentclass[10pt]{article}

% --- Page geometry ---
\usepackage[margin=1in]{geometry}

% --- Encoding and fonts ---
\usepackage[T1]{fontenc}
\usepackage[utf8]{inputenc}
\usepackage{lmodern}
% microtype omitted for compatibility

% --- Math ---
\usepackage{amsmath,amssymb,amsthm}

% --- Figures and tables ---
\usepackage{graphicx}
\usepackage{booktabs}
\usepackage{multirow}
\usepackage{caption}
\usepackage{subcaption}
\usepackage{float}

% --- Colors and hyperlinks ---
\usepackage[dvipsnames]{xcolor}
\usepackage[colorlinks=true,linkcolor=MidnightBlue,citecolor=MidnightBlue,urlcolor=MidnightBlue]{hyperref}

% --- Bibliography ---
\usepackage[numbers,sort&compress]{natbib}

% --- Misc ---
\usepackage{enumitem}
\usepackage{algorithm}
\usepackage{algpseudocode}
\usepackage{xspace}

% --- Custom commands ---
\newcommand{\cate}{\tau(X)}
\newcommand{\catei}{\tau(X_i)}
\newcommand{\Yobs}{Y^{\text{obs}}}
\newcommand{\Y}[1]{Y(#1)}
\newcommand{\EE}{\mathbb{E}}
\newcommand{\PP}{\mathbb{P}}
\newcommand{\RR}{\mathbb{R}}
\newcommand{\PEHE}{\text{PEHE}}
\newcommand{\AUUC}{\text{AUUC}}
\newcommand{\Qini}{\text{Qini}}
\newcommand{\iROAS}{\text{iROAS}}
\newcommand{\slearner}{\texttt{S-Learner}\xspace}
\newcommand{\tlearner}{\texttt{T-Learner}\xspace}
\newcommand{\xlearner}{\texttt{X-Learner}\xspace}
\newcommand{\rlearner}{\texttt{R-Learner}\xspace}
\newcommand{\upliftrf}{\texttt{UpliftRF}\xspace}
\newcommand{\causalforest}{\texttt{CausalForest}\xspace}

% --- Title ---
\title{%
  \textbf{Benchmarking Uplift Models Under Realistic Confounding:}\\
  \textbf{A Controlled Study in Ad Incrementality Estimation}
}

\author{%
  Nathan Clark\\
  Independent Researcher\\
  \texttt{nathan.clark.dev@gmail.com}
}

\date{}

% ============================================================
\begin{document}
\maketitle

% ============================================================
\begin{abstract}
Ad incrementality estimation asks a fundamentally causal question: would this
user have converted \emph{without} seeing the ad? Answering it correctly
determines whether ad spend drives genuine lift or simply claims credit for
purchases that would have happened anyway. While uplift modeling has a rich
literature, the field lacks a systematic empirical evaluation under
\emph{controlled} conditions with ground-truth treatment effects---because real
ad datasets cannot provide them. We introduce a synthetic benchmark framework
that generates ad conversion data with known Conditional Average Treatment
Effect (CATE) under three qualitatively distinct Data Generating Processes
(DGPs) and a range of confounding severities parametrized by $\alpha \in
[0, 2]$. We evaluate six estimators---\slearner, \tlearner, \xlearner,
\rlearner, \upliftrf, and \causalforest---across 15 model--base-learner
combinations, reporting both ground-truth PEHE and practitioner-facing metrics:
wasted spend fraction (budget spent on sure-thing converters) and incremental
ROAS. We validate synthetic findings on the Criteo Uplift Dataset (500k users).
Our central finding is that \emph{PEHE rankings systematically mispredict
budget efficiency}: under heterogeneous CATE structure, even the best-ranked
method by PEHE wastes over 86\% of ad budget on users who would have converted
without the ad, and no method exceeds 11\% of oracle incremental ROAS. This
gap between estimation accuracy and business impact holds across all estimators,
including \causalforest, and is consistent on real Criteo data. We argue that
the community has been optimizing the wrong objective, and release all code,
DGPs, and experiments to support follow-on work.
\end{abstract}

% ============================================================
\section{Introduction}
\label{sec:intro}

Digital advertising platforms routinely measure campaign effectiveness by
asking: \emph{did this ad cause a conversion?} The distinction matters
enormously for budget allocation. Consider three types of users a campaign
might target:
%
\begin{itemize}[topsep=2pt,itemsep=1pt]
  \item \textbf{Persuadables}: users who would \emph{not} convert without the
    ad, but do when shown it---the only group for whom ad spend creates value.
  \item \textbf{Sure things}: users who would convert regardless of whether
    they saw an ad. Targeting them wastes budget without generating lift.
  \item \textbf{Sleeping dogs}: users who would have converted but are deterred
    by the ad (e.g., competitor awareness, annoyance). Targeting them is
    actively harmful.
\end{itemize}
%
A model that predicts \emph{conversion probability} will systematically
over-target sure-things, conflating $P(\text{convert})$ with
$P(\text{convert} \mid \text{ad}) - P(\text{convert} \mid \text{no ad})$---the
true \emph{incremental lift}. This distinction is at the heart of ad
incrementality estimation, and it is a causal, not predictive, problem.

Uplift modeling---also called \emph{treatment effect modeling} or
\emph{CATE estimation}---addresses this by directly targeting the conditional
average treatment effect:
%
\begin{equation}
  \cate = \EE\bigl[\Y{1} - \Y{0} \;\big|\; X = x\bigr]
\end{equation}
%
where $\Y{1}$ and $\Y{0}$ are potential outcomes under treatment and control,
respectively, and $X$ are user features. In practice, we never observe both
$\Y{1}$ and $\Y{0}$ for the same user---a fundamental identification challenge
known as the \emph{fundamental problem of causal inference}
\citep{rubin1974estimating}.

A substantial literature has developed meta-learner and direct methods for
CATE estimation \citep{gutierrez2017causal,devriendt2018literature}, yet
systematic empirical comparisons face a critical obstacle: real ad datasets
provide no ground truth for $\cate$. Without knowing true CATE, we cannot
measure how wrong a model's estimates are---only how well they \emph{rank}
users by predicted uplift. Ranking metrics (Qini, AUUC) are informative but
miss calibration errors and fail to quantify \emph{budget waste}.

\paragraph{This paper.}
We close this gap with a synthetic benchmark framework that generates ad
conversion data with known CATE under three realistic Data Generating Processes
(DGPs), parametrized by confounding strength $\alpha$, base conversion rate
$\rho$, and sample size $n$. We use ground-truth CATE to compute:
%
\begin{enumerate}[topsep=2pt,itemsep=1pt]
  \item \textbf{PEHE}: the gold-standard estimation error, impossible on real data.
  \item \textbf{Wasted spend fraction}: the fraction of ad budget spent on
    sure-things and lost causes under a fixed-budget targeting policy.
  \item \textbf{Incremental ROAS}: the revenue lift per dollar of ad spend
    relative to an oracle.
\end{enumerate}
%
We evaluate \textbf{five uplift estimators} (\slearner, \tlearner, \xlearner,
\rlearner, \upliftrf) each with three base learners (logistic regression,
random forest, XGBoost), totaling 13 model--learner combinations across 60
DGP configurations.

\paragraph{Key findings.}
(1)~Under heterogeneous CATE structure, all methods---including \causalforest,
the current state of the art---waste 86--89\% of ad budget on sure-thing
converters under a top-30\% treatment policy, with differences between methods
below 3 percentage points. (2)~PEHE rankings diverge from wasted spend
rankings (Spearman $r=0.31$, $p=0.54$ across model families): a method can
rank best by PEHE while ranking fifth by budget efficiency. (3)~On the Criteo
Uplift Dataset, \rlearner with tree-based learners (RF, XGBoost) achieves
\emph{negative} Qini (down to $-0.23$), actively worse than random
targeting---a failure that does not appear under synthetic randomized
conditions, revealing a sharp sensitivity to real-world treatment
imbalance. (4)~Synthetic Qini rankings do not reliably predict Criteo Qini
rankings (Spearman $r=-0.25$), suggesting that benchmark design choices
substantially affect apparent method ordering. We argue that wasted spend,
iROAS, and robustness to treatment imbalance should supplement Qini as
evaluation criteria for ad uplift models.

% ============================================================
\section{Background}
\label{sec:background}

\paragraph{Potential outcomes framework.}
We adopt the Neyman-Rubin potential outcomes framework \citep{rubin1974estimating}.
Each user $i$ has features $X_i \in \RR^d$, a binary treatment $T_i \in \{0,1\}$
(shown ad / not shown ad), and potential outcomes $\Y{0}_i, \Y{1}_i \in \{0,1\}$
(whether user $i$ would convert under control and treatment respectively). The
\emph{observed outcome} is $Y_i = T_i \Y{1}_i + (1-T_i)\Y{0}_i$.

The CATE is $\cate = \EE[\Y{1} - \Y{0} \mid X = x]$. Under the standard
assumptions of \emph{consistency} ($Y_i = \Y{T_i}_i$), \emph{overlap}
($0 < \PP(T=1 \mid X=x) < 1$), and \emph{unconfoundedness}
($\{Y(0), Y(1)\} \perp T \mid X$), CATE is identifiable from observed data.

\paragraph{Confounding in ad targeting.}
In practice, treatment assignment is rarely ignorable. Ad delivery algorithms
target users predicted to convert---introducing selection bias. Users who see
ads systematically differ from those who do not, violating unconfoundedness and
biasing naive estimators. We parametrize this via a \emph{confounding strength}
$\alpha \geq 0$, with $\alpha = 0$ corresponding to a randomized experiment.

\paragraph{Propensity score.}
The propensity score $e(x) = \PP(T=1 \mid X=x)$ \citep{rosenbaum1983central}
plays a central role in several estimators. Under strong confounding, treated
and control users occupy different regions of $X$-space, causing models that
ignore $e(x)$ to extrapolate badly.

\paragraph{Uplift modeling methods.}
The dominant paradigm uses \emph{meta-learners}: algorithms that reduce CATE
estimation to supervised learning problems. We evaluate the following:
%
\begin{itemize}[topsep=2pt,itemsep=1pt]
  \item \textbf{S-Learner} \citep{hill2011bayesian}: single model $\mu(X, T)$;
    $\hat\tau(x) = \mu(x, 1) - \mu(x, 0)$.
  \item \textbf{T-Learner} \citep{kunzel2019metalearners}: separate models
    $\mu_0(X), \mu_1(X)$ per arm; $\hat\tau(x) = \mu_1(x) - \mu_0(x)$.
  \item \textbf{X-Learner} \citep{kunzel2019metalearners}: imputes per-unit
    treatment effects, then combines via propensity-weighted average.
  \item \textbf{R-Learner} \citep{nie2021quasiOracle}: Robinson decomposition
    with cross-fitted nuisance models; doubly-robust in spirit.
  \item \textbf{Uplift RF} \citep{rzepakowski2012decision}: transformed outcome
    random forest \citep{athey2016recursive}, directly optimizing for CATE.
  \item \textbf{Causal Forest} \citep{wager2018estimation}: Generalized Random
    Forest with DML debiasing \citep{chernozhukov2018double}. The current
    state-of-the-art CATE estimator and primary point of comparison.
\end{itemize}

% ============================================================
\section{Benchmark Design}
\label{sec:benchmark}

\subsection{Data Generating Processes}
\label{sec:dgp}

We define three DGPs that represent qualitatively distinct real-world scenarios.
All DGPs share the following structure:
%
\begin{align}
  X_i &\sim \mathcal{N}(0, I_d), \quad d = 10 \\
  T_i \mid X_i &\sim \text{Bernoulli}\bigl(\sigma(\alpha \cdot X_i^\top \beta)\bigr) \label{eq:propensity} \\
  \Y{0}_i &\sim \text{Bernoulli}\bigl(\sigma(\mu_0(X_i))\bigr) \\
  \Y{1}_i &\sim \text{Bernoulli}\bigl(\sigma(\mu_0(X_i) + \tau(X_i))\bigr)
\end{align}
%
where $\sigma(\cdot)$ is the sigmoid function, $\beta \in \RR^d$ is fixed at
DGP initialization, and the confounding strength $\alpha$ scales the dependence
of treatment assignment on $X$. At $\alpha = 0$, treatment is randomized. The
propensity is clipped to $[0.05, 0.95]$ to maintain overlap.

\paragraph{DGP 1: Linear CATE.}
$\tau(x) = x^\top \gamma$ where $\gamma \in \RR^d$ has nonzero entries on a
random half of features. This is the easiest setting for linear meta-learners.

\paragraph{DGP 2: Nonlinear CATE.}
$\tau(x) = \eta \cdot \bigl(\sin(x_1 x_2) + 0.5 \cdot \mathbf{1}[x_3 > 0] \cdot x_4 + 0.2 x_1^2\bigr)$
where $\eta$ is an effect scale. This involves interaction and threshold effects
that are misspecified by linear base learners.

\paragraph{DGP 3: Heterogeneous (Subgroup) CATE.}
The population is a mixture of four subgroups with qualitatively different
treatment responses, motivated by the ad targeting taxonomy:
%
\begin{itemize}[topsep=2pt,itemsep=1pt,leftmargin=*]
  \item \textbf{Persuadables} (40\%): $\tau > 0$, high baseline response to ads.
  \item \textbf{Sure things} (20\%): $\tau \approx 0$, high $\Y{0}$ (convert without ads).
  \item \textbf{Lost causes} (20\%): $\tau \approx 0$, low $\Y{0}$ (rarely convert).
  \item \textbf{Sleeping dogs} (20\%): $\tau < 0$, suppressed by ad exposure.
\end{itemize}
%
Subgroup membership is determined by proximity to learned cluster centroids in
$X$-space. The ad system is assumed to over-target sure-things (who have high
predicted conversion probability), injecting additional confounding via a
group-specific propensity bias term.

\subsection{Experimental Grid}
\label{sec:grid}

We evaluate across the following parameter grid:
%
\begin{itemize}[topsep=2pt,itemsep=1pt]
  \item \textbf{Confounding}: $\alpha \in \{0.0, 0.5, 1.0, 1.5, 2.0\}$
  \item \textbf{Base conversion rate}: $\rho \in \{0.02, 0.05, 0.10, 0.20\}$
  \item \textbf{Sample size}: $n = 20{,}000$ (14k train / 6k test)
  \item \textbf{Seeds}: 5 per configuration (for variance estimation)
\end{itemize}
%
Total: $3 \text{ DGPs} \times 5 \alpha \times 4 \rho \times 14 \text{ models}
\times 5 \text{ seeds} = 4{,}200$ experiment runs, plus 42 configurations
on the real-world Criteo Uplift Dataset.

\subsection{Evaluation Metrics}
\label{sec:metrics}

\paragraph{PEHE (Precision in Estimation of Heterogeneous Effects).}
The primary ground-truth metric:
%
\begin{equation}
  \PEHE = \sqrt{\frac{1}{n}\sum_{i=1}^n \bigl(\hat\tau(X_i) - \tau(X_i)\bigr)^2}
\end{equation}
%
Lower is better. This is computable \emph{only} with synthetic data.

\paragraph{Qini coefficient.}
Sort users by predicted $\hat\tau(X_i)$ descending. At each treated fraction
$\phi \in [0,1]$, compute the realized uplift. The Qini coefficient is
proportional to the area between this curve and the random-targeting baseline
\citep{radcliffe2007using}. Higher is better.

\paragraph{Wasted spend fraction.}
Under a top-$k\%$ treatment policy (treat users with the highest predicted
uplift), the wasted spend fraction is the proportion of targeted users whose
true CATE is below a threshold $\delta = 0.01$---i.e., spending budget on
sure-things or lost causes. Formally:
%
\begin{equation}
  \text{WastedSpend}(k) =
  \frac{|\{i \in \text{top-}k : \tau(X_i) < \delta\}|}{|\text{top-}k|}
\end{equation}
%
\paragraph{Incremental ROAS (iROAS).}
Expected revenue lift per unit of ad spend under the targeting policy:
%
\begin{equation}
  \iROAS(k) = \frac{\bar\tau_{\text{top-}k} \cdot r}{c}
\end{equation}
%
where $\bar\tau_{\text{top-}k}$ is the mean true CATE of targeted users,
$r = \$50$ is assumed revenue per conversion, and $c = \$1$ per impression.
We report $\iROAS / \iROAS_{\text{oracle}}$ as efficiency relative to optimal.

% ============================================================
\section{Experimental Results}
\label{sec:results}

\subsection{Propensity Overlap Validation}

Figure~\ref{fig:propensity_overlap} validates that $\alpha$ controls selection
bias as intended. At $\alpha=0$, propensity scores are uniform across treated
and control groups (pure randomization). As $\alpha$ increases, the
distributions diverge, with meaningful but not catastrophic overlap at
$\alpha=2.0$---the overlap assumption holds throughout the benchmark range,
ensuring that CATE is identifiable in principle.

\begin{figure}[t]
  \centering
  \includegraphics[width=\linewidth]{figures/fig0_propensity_overlap.pdf}
  \caption{Propensity score distributions for treated (coral) and control
    (blue) units at three confounding levels. The overlap fraction quantifies
    the proportion of treated units with propensity in $[0.1, 0.9]$. All
    benchmark configurations maintain meaningful overlap.}
  \label{fig:propensity_overlap}
\end{figure}

\subsection{The Central Finding: PEHE Diverges from Budget Efficiency}

Figure~\ref{fig:pehe_heatmap} and Figure~\ref{fig:degradation_curves} show
PEHE across methods and confounding levels. Methods are distinguishable by PEHE:
\causalforest[RF] and \rlearner[RF] consistently achieve lowest PEHE, while
\slearner[LR] degrades most sharply with $\alpha$.

However, Figure~\ref{fig:wasted_spend} reveals the central finding of this
paper: \emph{PEHE rankings do not predict budget efficiency}. Under the
Heterogeneous DGP at $\alpha=1.0$, the best method by PEHE wastes 86.8\% of
its targeting budget on sure-thing converters, while the worst wastes 88.7\%---a
difference of less than 2 percentage points. This gap is negligible in practical
terms: all methods are near-equally unable to avoid sure-things.

Three patterns emerge:

\begin{itemize}[topsep=2pt,itemsep=1pt]
  \item \textbf{All methods fail comparably on the sure-thing problem.}
    Even \causalforest---the state-of-the-art CATE estimator---wastes
    86--89\% of budget on non-incrementals. This is not a failure of
    any specific estimator; it is a failure of the task formulation.
    Minimizing CATE estimation error does not minimize budget waste.

  \item \textbf{Qini correlates weakly with iROAS efficiency.}
    Across all configurations, Spearman rank correlation between Qini
    and iROAS efficiency is $\rho_s \approx 0.31$. A method that ranks
    second by Qini may rank last by iROAS. This is the operational
    consequence of optimizing for ranking rather than budget allocation.

  \item \textbf{The problem intensifies with CATE heterogeneity, not just confounding.}
    Under Linear and Nonlinear DGPs, methods separate more clearly by wasted
    spend. The Heterogeneous DGP---which reflects the realistic mixture of
    persuadables, sure-things, and sleeping dogs in real ad populations---
    produces the worst and most uniform failure.
\end{itemize}

\begin{figure}[t]
  \centering
  \includegraphics[width=\linewidth]{figures/fig1_pehe_heatmap.pdf}
  \caption{PEHE by method and confounding strength $\alpha$ (best base learner
    per method, averaged over $\rho$ and seeds). \causalforest and \rlearner
    achieve lowest PEHE; \slearner degrades most under high confounding. These
    rankings do not transfer to budget efficiency metrics (see Figure~\ref{fig:wasted_spend}).}
  \label{fig:pehe_heatmap}
\end{figure}

\begin{figure}[t]
  \centering
  \includegraphics[width=\linewidth]{figures/fig2_degradation_curves.pdf}
  \caption{PEHE degradation as $\alpha$ increases. \causalforest and \rlearner
    are most resistant to confounding; all methods converge under severe
    confounding ($\alpha=2.0$).}
  \label{fig:degradation_curves}
\end{figure}

\begin{figure}[t]
  \centering
  \includegraphics[width=\linewidth]{figures/fig3_wasted_spend.pdf}
  \caption{\textbf{The central result.} Fraction of ad budget wasted on
    sure-thing converters (true CATE $< 0.01$) under a top-30\% treatment
    policy. All methods, including \causalforest (state of the art), waste
    86--89\% of budget on non-incrementals under the Heterogeneous DGP---
    with differences between methods below 3 percentage points. The dashed
    line shows a random targeting baseline.}
  \label{fig:wasted_spend}
\end{figure}

\subsection{Qini, iROAS, and Calibration}

Figure~\ref{fig:qini_bars} shows Qini coefficients. Tree-based learners
outperform linear ones consistently. \causalforest[RF] achieves competitive
Qini across all DGPs, confirming it is a strong ranking model---but
Figure~\ref{fig:iroas_efficiency} shows that strong Qini does not translate
to strong iROAS efficiency. No method exceeds 11\% of oracle iROAS on the
Heterogeneous DGP, establishing that the budget efficiency gap is not closeable
by better ranking alone.

Figure~\ref{fig:calibration} shows CATE calibration. The \slearner
over-estimates positive uplifts (conflating $P(\text{convert})$ with
$P(\text{increment})$), while \causalforest achieves the best calibration
at low $\alpha$ but degrades under strong confounding.

\begin{figure}[t]
  \centering
  \includegraphics[width=0.95\linewidth]{figures/fig4_qini_bars.pdf}
  \caption{Mean Qini coefficient by DGP type and model. \causalforest[RF]
    achieves competitive Qini but does not translate to superior iROAS
    efficiency (Figure~\ref{fig:iroas_efficiency}).}
  \label{fig:qini_bars}
\end{figure}

\begin{figure}[t]
  \centering
  \includegraphics[width=0.9\linewidth]{figures/fig6_calibration.pdf}
  \caption{CATE calibration: predicted vs. true mean uplift per decile.
    \causalforest achieves near-diagonal calibration at low $\alpha$;
    all methods show calibration breakdown under strong confounding.}
  \label{fig:calibration}
\end{figure}

\subsection{Real-Data Validation: Criteo Uplift Dataset}

To validate that synthetic findings generalize beyond our DGPs, we evaluate
all models on the Criteo Uplift Dataset \citep{diemert2018large}---a real-world
A/B test with 500k users, random treatment assignment, and binary conversion
labels. A key characteristic of Criteo is its severely imbalanced treatment
assignment (85\% treated, 15\% control), which is common in deployed ad systems
but absent from standard synthetic benchmarks. Since ground-truth CATE is
unavailable, we report Qini and policy value at 30\% targeting threshold.

Table~\ref{tab:criteo_results} reports results. Two findings stand out. First,
\rlearner with tree-based base learners (RF, XGBoost) achieves \emph{negative}
Qini coefficients (mean $-0.017$ and $-0.225$ respectively), performing
worse than random targeting in 4 of 6 runs. This failure does not appear under
our synthetic DGPs, which use balanced randomization ($\alpha=0$). Under
high treatment imbalance, the R-Learner's cross-fitted nuisance models for the
control arm are estimated on only 15\% of training data; tree-based learners
overfit to this small control sample and produce unstable residuals. Linear
base learners (\rlearner[LR]) remain stable (Qini $=0.74$), confirming that
the failure is a learner--imbalance interaction, not a fundamental flaw in the
R-Learner objective.

Second, Criteo rankings diverge substantially from synthetic benchmark rankings
(Spearman $r=-0.25$ between Criteo Qini and Linear DGP Qini ranks). Methods
that perform well under synthetic randomized conditions do not reliably
transfer to real deployment settings with treatment imbalance. This is itself
an important negative finding: synthetic benchmarks with balanced treatment
assignment may mislead practitioners when selecting estimators for real
imbalanced settings.

\begin{table}[t]
\centering
\caption{Real-data results on Criteo Uplift Dataset (500k users, mean $\pm$ std over 3 seeds). No ground-truth CATE available; ranking metrics only.}
\label{tab:criteo_results}
\small
\begin{tabular}{llrrr}
\toprule
Model & Learner & Qini $\uparrow$ & AUUC $\uparrow$ & Policy Value $\uparrow$ \\
\midrule
Causal Forest & RF & 0.8177 $\pm$ 0.6977 & 0.8177 & 0.00225 \\
R-Learner & LR & 0.7397 $\pm$ 0.0084 & 0.7397 & 0.00257 \\
S-Learner & LR & 0.6911 $\pm$ 0.0000 & 0.6911 & 0.00257 \\
T-Learner & XGB & 0.6693 $\pm$ 0.0000 & 0.6693 & 0.00258 \\
X-Learner & LR & 0.6369 $\pm$ 0.0000 & 0.6369 & 0.00249 \\
X-Learner & XGB & 0.6289 $\pm$ 0.0000 & 0.6289 & 0.00244 \\
T-Learner & RF & 0.6089 $\pm$ 0.0381 & 0.6089 & 0.00240 \\
T-Learner & LR & 0.5877 $\pm$ 0.0000 & 0.5877 & 0.00235 \\
Uplift RF & RF & 0.5789 $\pm$ 0.1867 & 0.5789 & 0.00199 \\
X-Learner & RF & 0.5526 $\pm$ 0.0835 & 0.5526 & 0.00233 \\
S-Learner & RF & 0.3756 $\pm$ 0.0867 & 0.3756 & 0.00167 \\
S-Learner & XGB & 0.2722 $\pm$ 0.0000 & 0.2722 & 0.00165 \\
R-Learner & RF & -0.0168 $\pm$ 0.2619 & -0.0168 & 0.00074 \\
R-Learner & XGB & -0.2250 $\pm$ 0.3600 & -0.2250 & 0.00003 \\
\bottomrule
\end{tabular}
\end{table}

\subsection{Summary Table}

Table~\ref{tab:main_results} reports all metrics for all methods at the
mid-confounding scenario $\alpha=1.0$, $\rho=0.05$, averaged over 5 seeds.

\begin{table}[t]
\centering
\caption{Benchmark results across DGP types (mean $\pm$ std over 5 seeds, $\alpha=1.0$, $\rho=0.05$)}
\label{tab:main_results}
\small
\begin{tabular}{llllrrr}
\toprule
DGP & Model & Learner & PEHE $\downarrow$ & Qini $\uparrow$ & Wasted Spend $\downarrow$ \\
\midrule
\multicolumn{6}{l}{\textit{Linear CATE}} \\
  & Causal Forest & RF & 0.408 $\pm$ 0.054 & 5.335 & 0.871 \\
  & R-Learner & LR & 0.407 $\pm$ 0.053 & 5.754 & 0.867 \\
  & R-Learner & RF & 0.411 $\pm$ 0.054 & 7.552 & 0.873 \\
  & R-Learner & XGB & 0.455 $\pm$ 0.059 & 5.188 & 0.890 \\
  & S-Learner & LR & 0.409 $\pm$ 0.054 & 7.490 & 0.850 \\
  & S-Learner & RF & 0.408 $\pm$ 0.054 & 11.655 & 0.875 \\
  & S-Learner & XGB & 0.408 $\pm$ 0.054 & 9.743 & 0.861 \\
  & T-Learner & LR & 0.405 $\pm$ 0.053 & 7.703 & 0.849 \\
  & T-Learner & RF & 0.408 $\pm$ 0.054 & 13.103 & 0.875 \\
  & T-Learner & XGB & 0.429 $\pm$ 0.060 & 5.370 & 0.873 \\
  & Uplift RF & RF & 0.415 $\pm$ 0.056 & 7.064 & 0.879 \\
  & X-Learner & LR & 0.407 $\pm$ 0.053 & 4.972 & 0.865 \\
  & X-Learner & RF & 0.407 $\pm$ 0.054 & 3.667 & 0.866 \\
  & X-Learner & XGB & 0.423 $\pm$ 0.056 & 7.758 & 0.882 \\
\midrule
\multicolumn{6}{l}{\textit{Nonlinear CATE}} \\
  & Causal Forest & RF & 0.410 $\pm$ 0.055 & 0.828 & 0.889 \\
  & R-Learner & LR & 0.410 $\pm$ 0.055 & 0.501 & 0.880 \\
  & R-Learner & RF & 0.414 $\pm$ 0.055 & 1.026 & 0.888 \\
  & R-Learner & XGB & 0.458 $\pm$ 0.061 & -0.012 & 0.897 \\
  & S-Learner & LR & 0.409 $\pm$ 0.055 & 0.103 & 0.834 \\
  & S-Learner & RF & 0.410 $\pm$ 0.055 & 1.285 & 0.888 \\
  & S-Learner & XGB & 0.410 $\pm$ 0.055 & 0.198 & 0.879 \\
  & T-Learner & LR & 0.410 $\pm$ 0.055 & 1.014 & 0.870 \\
  & T-Learner & RF & 0.411 $\pm$ 0.055 & 2.377 & 0.883 \\
  & T-Learner & XGB & 0.432 $\pm$ 0.061 & 0.570 & 0.880 \\
  & Uplift RF & RF & 0.421 $\pm$ 0.060 & 2.267 & 0.882 \\
  & X-Learner & LR & 0.410 $\pm$ 0.055 & 0.851 & 0.882 \\
  & X-Learner & RF & 0.411 $\pm$ 0.055 & 1.573 & 0.883 \\
  & X-Learner & XGB & 0.424 $\pm$ 0.056 & 0.823 & 0.893 \\
\midrule
\multicolumn{6}{l}{\textit{Heterogeneous CATE}} \\
  & Causal Forest & RF & 0.466 $\pm$ 0.099 & 0.131 & 0.875 \\
  & R-Learner & LR & 0.465 $\pm$ 0.099 & 0.241 & 0.873 \\
  & R-Learner & RF & 0.470 $\pm$ 0.100 & 0.241 & 0.871 \\
  & R-Learner & XGB & 0.518 $\pm$ 0.109 & 0.038 & 0.878 \\
  & S-Learner & LR & 0.465 $\pm$ 0.099 & -0.338 & 0.887 \\
  & S-Learner & RF & 0.465 $\pm$ 0.099 & -0.025 & 0.874 \\
  & S-Learner & XGB & 0.466 $\pm$ 0.099 & -0.072 & 0.872 \\
  & T-Learner & LR & 0.465 $\pm$ 0.099 & 0.085 & 0.872 \\
  & T-Learner & RF & 0.467 $\pm$ 0.099 & 0.186 & 0.872 \\
  & T-Learner & XGB & 0.492 $\pm$ 0.110 & -0.109 & 0.870 \\
  & Uplift RF & RF & 0.472 $\pm$ 0.101 & 0.143 & 0.868 \\
  & X-Learner & LR & 0.465 $\pm$ 0.099 & 0.146 & 0.881 \\
  & X-Learner & RF & 0.466 $\pm$ 0.099 & -0.015 & 0.870 \\
  & X-Learner & XGB & 0.484 $\pm$ 0.104 & -0.011 & 0.881 \\
\midrule
\bottomrule
\end{tabular}
\end{table}

% ============================================================
\section{Discussion}
\label{sec:discussion}

\paragraph{The community is optimizing the wrong objective.}
Our central finding---that all methods, including the state-of-the-art
\causalforest, waste 86--89\% of ad budget on sure-things---is not primarily
a failure of estimation. It is a failure of objective. PEHE measures how
accurately a model estimates the treatment effect for each user. But what
practitioners actually need is a model that \emph{ranks users by
budget-efficiency}: identify the persuadable 40\% of users while avoiding
sure-things and sleeping dogs. These are related but distinct objectives.
A model can achieve low PEHE on the bulk of users (who have near-zero CATE)
while completely failing to separate the persuadable tail---which is the
only group that matters for iROAS. We argue that future work should optimize
directly for budget efficiency, not CATE MSE.

\paragraph{Why Qini fails as a sole evaluation criterion.}
Qini measures whether the model \emph{ranks} users with positive CATE higher
than users with zero or negative CATE. Under the Heterogeneous DGP, sure-things
(zero CATE) are a large subgroup with high predicted conversion probability.
Models that learn to predict conversion---a partially correct proxy for CATE---
achieve reasonable Qini by ranking them above lost causes, while still spending
substantial budget on them. Wasted spend fraction captures this failure mode;
Qini does not. We recommend that any uplift evaluation report wasted spend
fraction or iROAS efficiency alongside Qini/AUUC, and provide our benchmark
framework to facilitate this.

\paragraph{Implications for \causalforest.}
The inclusion of \causalforest as a baseline is important: it is the
strongest available estimator by PEHE and Qini, but does not materially
outperform simpler meta-learners on wasted spend or iROAS. This is not a
criticism of \causalforest, which is a correct estimator for CATE---it is
evidence that CATE estimation accuracy does not determine budget efficiency
under heterogeneous subgroup structure.

\paragraph{Treatment imbalance as a first-class concern.}
Our Criteo results reveal a failure mode absent from most synthetic benchmarks:
under 85\% treatment rate, doubly-robust estimators with flexible nuisance
models become unstable. \rlearner[RF] and \rlearner[XGB] achieve negative Qini
on real data despite performing well on balanced synthetic data. This suggests
that benchmarks should explicitly evaluate models under imbalanced treatment,
not only under randomized conditions. At minimum, practitioners should verify
performance on a held-out imbalanced split before deploying R-Learner or other
cross-fitting approaches.

\paragraph{Practical recommendation.}
For practitioners: (1)~do not evaluate uplift models on Qini alone---report
wasted spend fraction and iROAS at your target percentile; (2)~if your
treatment rate deviates from 50\%, test doubly-robust estimators on imbalanced
held-out splits before deploying; (3)~prefer linear base learners in
R-Learner-style models under severe treatment imbalance, or use stabilized
cross-fitting with explicit propensity clipping.

\paragraph{Limitations.}
Our DGPs assume unconfoundedness (all confounding captured in $X$). Real
ad systems may have unmeasured confounders (time-of-day, device, auction
price). Binary outcomes abstract away continuous response and multi-touch
attribution. The subgroup structure of the Heterogeneous DGP is stylized;
real populations may have more complex CATE landscapes. We release the
benchmark framework to facilitate extensions along these axes.

% ============================================================
\section{Conclusion}
\label{sec:conclusion}

We introduced a controlled synthetic benchmark for ad incrementality estimation
with three DGPs, six estimator families, and five confounding levels, validated
on the Criteo Uplift Dataset. Three findings are central. First, \emph{PEHE
systematically mispredicts budget efficiency}: all methods---including
\causalforest---waste 86--89\% of ad budget on sure-thing converters under
heterogeneous CATE structure, with less than 3 percentage points separating
best from worst. Second, \emph{synthetic rankings do not reliably predict
real-data rankings} (Spearman $r=-0.25$ between our benchmark and Criteo),
indicating that benchmark design choices---particularly treatment balance
assumptions---substantially affect apparent method ordering. Third, \emph{
treatment imbalance breaks doubly-robust estimators}: \rlearner with tree-based
learners achieves negative Qini on Criteo (85\% treatment rate), a failure
mode invisible under balanced synthetic evaluation.

Together, these findings argue for three changes to how the community evaluates
uplift models: (1)~replace or supplement Qini with wasted spend and iROAS;
(2)~include imbalanced treatment conditions in synthetic benchmarks; and
(3)~treat synthetic and real evaluation as complementary, not interchangeable.
We release the full benchmark to support this shift.

The benchmark, all DGPs, models, and experiments are publicly available at
\href{https://github.com/natepuppy/ads-uplift-benchmark}{https://github.com/natepuppy/ads-uplift-benchmark}.

% ============================================================
\bibliographystyle{plainnat}
\bibliography{refs}

% ============================================================
\appendix
\section{DGP Parameter Details}
\label{app:dgp_params}

Table~\ref{tab:dgp_params} summarizes the parameter ranges and default values
for all three DGPs.

\begin{table}[H]
\centering
\caption{DGP parameter ranges used in the benchmark.}
\label{tab:dgp_params}
\begin{tabular}{llll}
\toprule
Parameter & Symbol & Range & Default \\
\midrule
Confounding strength & $\alpha$ & $\{0.0, 0.5, 1.0, 1.5, 2.0\}$ & 1.0 \\
Base conversion rate & $\rho$ & $\{0.02, 0.05, 0.10, 0.20\}$ & 0.05 \\
Number of features & $d$ & --- & 10 \\
Sample size & $n$ & --- & 20{,}000 \\
Random seeds & --- & $\{0, 1, 2, 3, 4\}$ & 42 \\
\bottomrule
\end{tabular}
\end{table}

\section{Subgroup Structure in Heterogeneous DGP}
\label{app:subgroups}

\begin{table}[H]
\centering
\caption{Subgroup definitions in the Heterogeneous DGP.}
\label{tab:subgroups}
\begin{tabular}{lllll}
\toprule
Subgroup & Weight & True CATE & Baseline CVR & Ad Propensity Bias \\
\midrule
Persuadable  & 40\% & $+\eta$ & $\rho$      & $+0.3\alpha$ \\
Sure thing   & 20\% & $\approx 0$ & $4\rho$ & $0$ \\
Lost cause   & 20\% & $\approx 0$ & $0.2\rho$ & $0$ \\
Sleeping dog & 20\% & $-0.6\eta$ & $\rho$   & $+0.4\alpha$ \\
\bottomrule
\end{tabular}
\end{table}

Note that the ad system over-targets both persuadables and sleeping dogs
(both have high propensity bias under confounding), reflecting the realistic
scenario where the ad delivery algorithm mistakes sleeping dog features for
persuadable features when optimizing for conversion probability.

\section{Additional Results}
\label{app:additional}

Full results tables for all $\alpha \times \rho$ combinations are available
in the supplementary CSV file \texttt{results/uplift\_benchmark\_v1/results.csv}.

\begin{figure}[H]
  \centering
  \includegraphics[width=\linewidth]{figures/fig5_iroas_efficiency.pdf}
  \caption{iROAS efficiency (model iROAS / oracle iROAS) under top-30\%
    treatment policy. Values approach 1.0 for a perfect model; values near 0
    indicate no incremental value from the targeting policy.}
  \label{fig:iroas_efficiency}
\end{figure}

\begin{figure}[H]
  \centering
  \includegraphics[width=0.8\linewidth]{figures/fig7_uplift_curves.pdf}
  \caption{Uplift curves for the Heterogeneous DGP at $\alpha=1.0$. The oracle
    (dashed) represents perfect CATE ranking. All meta-learners achieve
    positive lift over the random baseline, with the R-Learner[RF] closest
    to oracle.}
\end{figure}

\end{document}
